% ============================================================
% CONCLUSION SECTION
% ============================================================

\section{Conclusion}\label{sec:conclusion}

This study demonstrates that offline RL algorithm selection profoundly impacts policy interpretability independent of performance in the gym-sepsis simulator. Our systematic LEG-based benchmark shows that Conservative Q-Learning achieves substantially stronger saliency magnitude than vanilla DQN, spanning multiple orders of magnitude, while achieving comparable survival outcomes. While all algorithms achieve similar overall survival rates, CQL substantially outperforms DQN on high-severity patients with elevated SOFA scores, representing a clinically meaningful performance gap. CQL's strong, clinically coherent saliency patterns emphasizing blood pressure and lactate align with medical guidelines and provide quantitative feature importance suitable for regulatory review.

Our findings within this single simulator suggest that conservative offline RL methods warrant consideration for safety-critical healthcare applications where interpretability is essential. CQL's conservative penalty term biases the Q-function toward values supported by the training dataset, inheriting the behavioral policy's interpretable structure based on threshold-based heuristics. This interpretability-by-design mechanism offers an alternative to post-hoc explainability methods. However, generalization beyond the gym-sepsis environment remains uncertain. Our conclusions are based on a single sepsis treatment simulator with known limitations including high baseline survival and learned dynamics, a limited set of representative algorithms, and LEG as the sole interpretability metric. Claims about broader applicability to other clinical domains such as autonomous driving or financial trading or real-world deployment require prospective validation through multi-center clinical trials, human-in-the-loop studies with clinicians, and evaluation on real EHR data. We establish a rigorous evaluation framework combining performance metrics, SOFA-stratified analysis, and quantitative interpretability assessment through LEG saliency with clinical coherence, providing a template for future healthcare RL benchmarking studies.

Future work should pursue prospective clinical validation, integrate domain knowledge into CQL training through feature-aware regularization, develop healthcare-specific interpretability metrics including stability, parsimony, and actionability, and conduct human-in-the-loop experiments to establish causality between interpretability and clinical decision quality. By prioritizing interpretability alongside performance in algorithm design, we can develop clinical AI systems that earn the trust and adoption of clinicians while improving patient care.

% End of Conclusion section